\documentclass[12pt]{article}%
\usepackage{amsfonts}
\usepackage{fancyhdr}
\usepackage{comment}
\usepackage[a4paper, top=2.5cm, bottom=2.5cm, left=2.2cm, right=2.2cm]%
{geometry}
\usepackage{times}
\usepackage{amsmath}
\usepackage{changepage}
\usepackage{amssymb}
\usepackage{graphicx}

\setcounter{MaxMatrixCols}{30}
\newtheorem{theorem}{Theorem}
\newtheorem{acknowledgement}[theorem]{Acknowledgement}
\newtheorem{algorithm}[theorem]{Algorithm}
\newtheorem{axiom}{Axiom}
\newtheorem{case}[theorem]{Case}
\newtheorem{claim}[theorem]{Claim}
\newtheorem{conclusion}[theorem]{Conclusion}
\newtheorem{condition}[theorem]{Condition}
\newtheorem{conjecture}[theorem]{Conjecture}
\newtheorem{corollary}[theorem]{Corollary}
\newtheorem{criterion}[theorem]{Criterion}
\newtheorem{definition}[theorem]{Definition}
\newtheorem{example}[theorem]{Example}
\newtheorem{exercise}[theorem]{Exercise}
\newtheorem{lemma}[theorem]{Lemma}
\newtheorem{notation}[theorem]{Notation}
\newtheorem{problem}[theorem]{Problem}
\newtheorem{proposition}[theorem]{Proposition}
\newtheorem{remark}[theorem]{Remark}
\newtheorem{solution}[theorem]{Solution}
\newtheorem{summary}[theorem]{Summary}
\newenvironment{proof}[1][Proof]{\textbf{#1.} }{\ \rule{0.5em}{0.5em}}

\newcommand{\Q}{\mathbb{Q}}
\newcommand{\R}{\mathbb{R}}
\newcommand{\C}{\mathbb{C}}
\newcommand{\Z}{\mathbb{Z}}

%%%%
% ME %
%%%%

\usepackage{listings}
\usepackage{color}
\usepackage{xcolor}
\usepackage{mdframed}

\definecolor{dkgreen}{rgb}{0,0.6,0}
\definecolor{gray}{rgb}{0.5,0.5,0.5}
\definecolor{mauve}{rgb}{0.58,0,0.82}

\lstset{%frame=tb,
language=Matlab,
aboveskip=3mm,
belowskip=3mm,
showstringspaces=false,
columns=flexible,
basicstyle={\small\ttfamily},
numbers=none,
numberstyle=\tiny\color{gray},
keywordstyle=\color{blue},
commentstyle=\color{dkgreen},
stringstyle=\color{mauve},
breaklines=true,
breakatwhitespace=true,
tabsize=3
}

\usepackage{outlines}
\usepackage{enumitem}
%\setenumerate[1]{label=\Roman*.}
%\setenumerate[2]{label=\Alph*.}
%\setenumerate[3]{label=\roman*.}
%\setenumerate[4]{label=\alph*.}

\usepackage{titlesec} 

\titleformat{\subsection}[runin]
  {\normalfont\large\bfseries}{\thesubsection}{1em}{}
\titleformat{\subsubsection}[runin]
  {\normalfont\normalsize\bfseries}{\thesubsubsection}{1em}{}

% for enumerate spacing
\newenvironment{tight_enum}{
\begin{enumerate}[label=\alph*.]
\setlength{\itemsep}{0pt}
\setlength{\parskip}{0pt}
}{\end{enumerate}}

\usepackage{exercise}

\setlength\parindent{0pt}

\newcommand{\code}[1]{\texttt{#1}}
\newcommand{\shp}{$>>>$ }
\newcommand{\tab}{\hspace*{45pt}}

\usepackage{multicol}
\setlength{\columnsep}{0.6cm}

\begin{document}

\title{Homework 04: Chapter 05}

\author{EEC 3150}
\date{}

\maketitle

\begin{center}
\fbox{\fbox{\parbox{5.5in}{\centering
For each Python exercise, write a separate Python script. Submit all scripts on Blackboard under the appropriate assignment. Name your script files with the following format: \\
\vspace{10pt}
HW$<$assignment number$>$\_Ex$<$exercise number$>$\_$<$FirstInitial$><$LastName$>$.py \\
\vspace{10pt}
Example: HW01\_Ex01\_GWest.py \\
\vspace{10pt}
For short answer or math exercises, either do your work on a sheet of paper and upload a picture of it to Blackboard or submit some kind of text file or Word document to Blackboard. You can place all such exercises in a single document as long as you clearly label each exercise. These files should follow a similar naming convention as the Python scripts.
}}}
\end{center}


\pagebreak


\begin{multicols*}{2}

%%% NEW EXERCISE %%%
\begin{Exercise}[origin={10 points, Python}]

\begin{enumerate}[leftmargin=*]
\item Write a program which will read the file CSVData.csv and store its contents as a 2D list of floats. For example, if the file looked like this

\code{1,2,3\\ %
4,5,6\\ %
7,8,9
}

then your 2D array should look like this

\code{[[1,2,3],\\ \hspace{0pt}
[4,5,6],\\ \hspace{0pt}
[7,8,9]]
}
\item Print the 2D array with the following code:

\code{for x in data: \\
\tab print(x)
}

(assuming \code{data} is the name of your 2D list).
\end{enumerate}

\end{Exercise}

%%% NEW EXERCISE %%%
\begin{Exercise}[origin={10 points, Python}]

\begin{enumerate}[leftmargin=*]
\item Write a program (not a function) that will read the files
\begin{itemize}
	\setlength\itemsep{0pt}
	\item Mercury.txt
	\item Venus.txt
	\item Earth.txt
	\item Mars.txt
	\item Jupiter.txt
\end{itemize}
and store the radius of each planet in a float variable. NOTE: in the first four files, the radius is called \code{Mean radius (km)} while in the fifth file it is called \code{Equat. radius (1 bar)}.
\item Print the radii of each planet and label them clearly.
\end{enumerate}

\end{Exercise}

\vfill\null
\columnbreak

%%% NEW EXERCISE %%%
\begin{Exercise}[origin={10 points, Python}]

\begin{enumerate}[leftmargin=*]
	\item Write a program (not a function) which will  read the file ``bible-psalms.txt'' and create a nested dictionary \code{psalms}. The keys of this dictionary will be the chapter numbers and the values will be another dictionary. The keys of these inner dictionaries will be a verse number and the values will be a string that is the verse's text. NOTE: the tab (\textbackslash  t) escape character is used in each line of the file to separate the verse number and the verse text.
	\item Give a demo of indexing a particluar verse from the dictionary. Here is an example (you may choose your preferred chapter/verse):
	
	\code{chapter = 1 \\
	verse = 1 \\
	print(psalms[chapter][verse])
	}
	
	The output for the above example should look like:
	
	\code{Blessed is the man that walketh not in the counsel of the ungodly, nor standeth in the way of sinners, nor sitteth in the seat of the scornful.
	}	
\end{enumerate}

\end{Exercise}

\vfill\null
\end{multicols*}




\end{document}






